\section{Introduction}
\label{sec:intro}

Video-to-audio (V2A) synthesis asks a model to produce an audio
waveform that is both semantically plausible and temporally
synchronized with a silent video. The current state of the art,
MMAudio \citep{cheng2024mmaudio}, frames V2A as conditional flow
matching \citep{lipman2022flow}: an MMDiT
\citep{esser2024scaling} transformer denoises a VAE-encoded
mel-spectrogram latent given three conditioning streams --- a
CLIP \citep{radford2021learning} semantic embedding, a
Synchformer \citep{iashin2024synchformer} temporal track, and
optional CLIP text --- and the model is trained jointly on
audio--visual and audio--text corpora at a roughly $1{:}1$
sampling ratio. MMAudio reports new state-of-the-art
FD$_{\text{PaSST}}$, IB-score, and DeSync with only 157\,M
parameters, by exploiting the much larger pool of
audio--text data (AudioCaps \citep{kim2019audiocaps},
Clotho, WavCaps \citep{mei2024wavcaps}) on top of VGGSound
\citep{chen2024vggsound}.

\subsection{MMAudio's objective is purely pointwise}
\label{sec:intro:pointwise}
A careful reading of \citet{cheng2024mmaudio} shows that
\emph{every} cross-modal coupling in MMAudio is
instance-level. The training loss is exactly the conditional
flow-matching loss
\begin{equation}
\label{eq:mmaudio-fm}
  \E_{t, q(x_0), q(x_1, \mathbf{C})}
    \bigl\lVert v_\theta(t, \mathbf{C}, x_t)
      - (x_1 - x_0) \bigr\rVert_2^2,
\end{equation}
with no auxiliary contrastive, alignment, or distribution-matching
term. The cross-modal machinery enters only through the
architecture:
\begin{itemize}[leftmargin=*,topsep=2pt]
  \item \textbf{Joint MMDiT attention.} Each of the $N_1{=}4$
    multimodal blocks concatenates the (Q, K, V) triples from
    audio, CLIP, and text streams and applies one joint scaled
    dot-product attention. The inner product $q_a \cdot k_v$
    binds a single audio token to a single video token.
  \item \textbf{Conditional Synchronization Module.}
    Synchformer features at 24\,fps are upsampled to the
    31.25\,fps audio latent rate and injected through a
    frame-aligned adaLN, whose scale/shift parameters modulate
    each audio token by the temporally co-located video
    feature.
  \item \textbf{Empty-token masking.} For audio-only samples,
    visual streams are replaced by learnable
    $\varnothing_v, \varnothing_{\mathrm{syn}}$ tokens, preserving
    the same pointwise plumbing.
\end{itemize}
In the terminology of \citet{wang2020understanding}, every
supervision signal in MMAudio controls \emph{alignment}
--- whether paired video and audio features end up near each
other --- but nothing in its objective controls
\emph{uniformity} or, more specifically, the
\emph{pairwise} geometry of either modality. If clips $i$ and
$j$ look similar on the video side, no term in
Eq.\,(\ref{eq:mmaudio-fm}) encourages the corresponding audio
latents to be similar on the audio side. We call this the
\emph{pointwise--relational gap}.

\subsection{The pointwise--relational gap matters}
Pointwise losses provably control the average distance between
paired positives but leave the marginal geometries free
\citep{wang2020understanding}. For V2A this is a real concern:
a listener's judgment of plausibility is not just ``does this
audio match this video?'' but also ``do audios of similar
videos sound similar?'' --- a relational property of the
distribution, not of a single pair. The standard remedy in
retrieval settings is a contrastive negatives term, which adds
a uniformity-on-the-hypersphere objective. That term is still
pointwise per pair; it produces a repulsion but not an explicit
geometric match between $\DV$ and $\DA$.

\subsection{Empirical motivation: MMAudio's geometry is misaligned}
\label{sec:intro:motivation}
The previous subsection gave a loss-level argument; we also want
an empirical one. Using the released MMAudio \texttt{small\_16k}
checkpoint with \emph{no additional training}, we extract
per-clip video and audio representations on the VGGSound test
split and measure four quantities that probe relational
alignment. The script is
\texttt{experiments/diagnose\_geometry.py} and a one-GPU Slurm
wrapper is \texttt{jobs/diagnose\_geometry.job}. Each probe is
computed for all four GW representation variants
(Sec.\,\ref{sec:design}).

\begin{enumerate}[leftmargin=*,topsep=2pt]
  \item \textbf{Pairwise-distance scatter.} For a batch of
    $B{\approx}2000$ test clips, compute the
    $B \times B$ squared-distance matrices $\DV$ and $\DA$ and
    scatter $\DV[i,k]$ against $\DA[i,k]$ on the upper
    triangle. An isometric embedding would be a tight
    $45\degree$ line; we report the Pearson and Spearman
    correlations and expect both to be far below $1$
    (Fig.\,\ref{fig:motivation}\,(a)).
  \item \textbf{kNN-graph agreement.} For
    $k \in \{5, 10, 20\}$, build the video-side and audio-side
    $k$-nearest-neighbor graphs and report the mean per-sample
    Jaccard overlap of neighborhoods
    (Fig.\,\ref{fig:motivation}\,(b)). A purely pointwise model
    can place each positive pair near each other while scrambling
    the neighborhood around it; this probe reads out that
    scramble directly.
  \item \textbf{Class-pair rank correlation.} Using VGGSound
    labels, compute class centroids on each side, form the
    between-class distance matrices $\bar\DV, \bar\DA$, and report
    Kendall's $\tau$ between their upper triangles. Low $\tau$
    means ``visually similar classes'' are not ``audibly similar
    classes'' under the shipped checkpoint.
  \item \textbf{Raw GW / FGW at init.} Evaluate the entropic
    GW loss of Sec.\,\ref{sec:method:gw} on random
    $B{=}64$ batches from the test reps. A large finite value
    is quantitative evidence that $\DV$ and $\DA$ are far apart
    in GW geometry \emph{before} any relational regularization
    is applied.
\end{enumerate}

We also render the optimal coupling $\Tstar$ from
\texttt{compute\_gw\_regularization} for the top-20
most-frequent test classes
(Fig.\,\ref{fig:motivation}\,(c)) and report the diagonal-mass
fraction. Under a relationally aligned model the heatmap is
strongly block-diagonal; under a pointwise-only one, mass
leaks off-diagonal and the diagonal-mass fraction is close to
the uniform baseline $1/n_{\text{class}}$.

\begin{figure}[t]
\centering
\includegraphics[width=0.32\linewidth]{motivation/scatter_global.png}\hfill
\includegraphics[width=0.32\linewidth]{motivation/knn.png}\hfill
\includegraphics[width=0.32\linewidth]{motivation/coupling.png}
\caption{Motivating diagnostics on the released MMAudio
\texttt{small\_16k} checkpoint (no fine-tuning). \textbf{(a)}
Pairwise-distance scatter $\DV$ vs $\DA$ on 2000 VGGSound-test
clips (\texttt{global} variant); Pearson and Spearman
correlations in title. \textbf{(b)} Mean Jaccard overlap of
video- and audio-side $k$-NN graphs as $k$ varies, across the
four GW variants. \textbf{(c)} Class-pair GW coupling heatmap
$\Tstar[i,j]$ over the 20 most frequent test classes; diagonal
mass fraction annotated. All three panels support the claim that
MMAudio's pointwise-only objective leaves the two modalities
relationally misaligned. Produced by
\texttt{jobs/diagnose\_geometry.job}; CSV and LaTeX tables are
written to \texttt{analysis/motivation/metrics.\{csv,tex\}}.}
\label{fig:motivation}
\end{figure}

\begin{table}[t]
\centering\small
\begin{tabular}{@{}lcccccc@{}}
\toprule
variant & $r_\text{Pearson}(\DV,\DA)$ & $\text{Jaccard}_{k=10}$ &
$\tau_\text{class}$ & $\LGW^{(\text{init})}$ &
$\text{erank}(V)$ & $\text{erank}(A)$ \\
\midrule
\texttt{global}    & --- & --- & --- & --- & --- & --- \\
\texttt{projected} & --- & --- & --- & --- & --- & --- \\
\texttt{c\_g}      & --- & --- & --- & --- & --- & --- \\
\texttt{fused}     & --- & --- & --- & --- & --- & --- \\
\bottomrule
\end{tabular}
\caption{Relational-alignment diagnostics on the released MMAudio
checkpoint. Filled from
\texttt{analysis/motivation/metrics.csv} after running
\texttt{jobs/diagnose\_geometry.job}. The hypothesis:
$r_\text{Pearson}$, $\text{Jaccard}$, $\tau_\text{class}$ are all
small (bounded away from $1$) and $\LGW^{(\text{init})}$ is large,
demonstrating the gap the regularizer of
Sec.\,\ref{sec:method:reg} targets.}
\label{tab:motivation}
\end{table}

\subsection{Our contribution}
\label{sec:intro:contribution}
We close the pointwise--relational gap by adding exactly one
term to the MMAudio training objective: an entropic
Gromov--Wasserstein (GW) \citep{peyre2016gromov} penalty between
the intra-video and intra-audio pairwise squared-distance
matrices of each mini-batch. To our knowledge this is the first
relational-geometry regularizer for a V2A flow-matching model;
it is orthogonal to, and composes with, every pointwise
mechanism MMAudio already uses. Concretely:
\begin{enumerate}[leftmargin=*,topsep=2pt]
  \item \textbf{Formalization.} We derive the regularized
    objective $\mathcal{L} = \LFM + \lambda_{\mathrm{GW}} \LGW$
    and show that the GW term is null-space-orthogonal to the
    flow-matching term on the individual-example axis: $\LGW$
    depends only on mean-pooled per-sample representations, so it
    cannot be satisfied by shifting single predictions.
  \item \textbf{Implementation.} We solve entropic GW with
    Peyr\'e--Sinkhorn mirror descent
    \citep{peyre2016gromov,cuturi2013sinkhorn} in float32 inside
    an \texttt{autocast(enabled=False)} block, preserving the
    rest of the network's bfloat16 training. The GW code is 232
    lines of \texttt{mmaudio/model/gw\_regularization.py} and
    does not touch the network architecture.
  \item \textbf{Variants and knobs.} We enumerate five design
    axes D1--D5: which representation pair to compare
    (\texttt{global / projected / c\_g / fused}), the penalty
    weight $\lambda_{\mathrm{GW}}$, detach direction, schedule,
    and interaction with Synchformer.
    \texttt{fused} additionally uses FGW \citep{vayer2020fused}
    to blend relational and pointwise costs.
  \item \textbf{Evaluation.} We evaluate on MMAudio's own
    VGGSound benchmark with its four-metric suite
    (FD$_{\text{PaSST}}$, IS, IB-score, DeSync via
    \texttt{av-benchmark}), plus a 30-class held-out OOD split
    (\texttt{experiments/make\_ood\_split.py}) that probes
    whether the relational term helps on genuinely unseen
    classes, and coupling-matrix diagnostics
    (\texttt{experiments/analyze\_couplings.py}) that read
    out the learned $\Tstar$ directly.
\end{enumerate}

\subsection{Scope of the contribution}
We modify only the loss. The MMDiT stack, CLIP encoder,
Synchformer branch, Conditional Synchronization Module, VAE,
and BigVGAN \citep{lee2023bigvgan} vocoder are exactly those of
MMAudio. We run at the public \texttt{small\_16k}
configuration (157\,M parameters, 25 Euler steps, CFG 4.5, Post-Hoc
EMA) and do not retune any of MMAudio's other hyperparameters.
This makes the delta fully attributable to the added term:
any change in metrics, geometry, or coupling structure is a
consequence of $\lambda_{\mathrm{GW}} \LGW$, not of a confounding
architectural modification.

\subsection{Thesis roadmap}
Sec.\,\ref{sec:background} recaps flow matching, GW / FGW, and
the specific MMAudio mechanisms our term sits alongside.
Sec.\,\ref{sec:method} formalizes the regularizer.
Sec.\,\ref{sec:design} lists the design axes.
Sec.\,\ref{sec:experiments} states H1--H5 and the ablation grid.
Sec.\,\ref{sec:results} is filled after runs complete.
Sec.\,\ref{sec:conclusion} positions the work against MMAudio's
own acknowledged limitations and names extensions.
